%========================================================================
%  SOUTENANCE DE STAGE --- Nathan Legendre --- INSA Rouen Normandie
%  Ecocea La Fabrique --- Squelette Beamer (theme Simple)
%  >>> COMPILER DEPUIS CE DOSSIER (presentation/) : latexmk -pdf presentation.tex
%========================================================================

\documentclass[aspectratio=169,xcolor=dvipsnames]{beamer}

\makeatletter
\def\input@path{{packages/}}
\makeatother

\usetheme{Simple}

\usepackage[T1]{fontenc}
% Typo française : ton install (rapport) a déjà le paquet. Décommente la ligne suivante :
% \usepackage[french]{babel}
\usepackage{hyperref}
\usepackage{graphicx}
\usepackage{booktabs}

% --- Placeholder visuel : a remplacer par tes captures / figures du rapport ---
\newcommand{\placeholder}[2]{%
  \begin{center}\setlength{\fboxsep}{0pt}%
  \colorbox{InvisibleBlue}{\framebox[0.82\linewidth]{\rule{0pt}{#1}%
  \textcolor{gray}{\footnotesize [#2]}}}\end{center}}

%------------------------------------------------------------------------
\title[Démarche QA/IA --- Polène]{Industrialiser la QA d'un e-commerce multi-pays}
\subtitle{Une démarche ingénieur appuyée sur l'IA générative --- Stage chez Ecocea La Fabrique}
\author[Nathan Legendre]{Nathan Legendre}
\institute[INSA Rouen]{%
  Département Informatique \\
  INSA Rouen Normandie
  \vskip 3pt}
\date{\today}

%------------------------------------------------------------------------
\begin{document}

\begin{frame}
  \titlepage
\end{frame}

\begin{frame}{Plan}
  \tableofcontents
\end{frame}

% ===== CONTENU =====
%========================================================================
\section{Contexte \& fil rouge}
% ⏱ ~4 min  |  Objectif : poser le décor, montrer l'étendue du rôle, ENTONNOIR vers la question phare
%========================================================================

\begin{frame}{Ecocea La Fabrique \& le projet Polène}
  % 🎤 ~1 min. Qui est l'agence, sur quoi j'étais. Rester factuel et rapide.
  \begin{columns}[T]
    \column{.55\textwidth}
    \begin{itemize}
      \item Intégrateur e-commerce ($\sim$60 ingénieurs), écosystème \alert{Shopify} \& SFCC
      \item Comptes retail haut de gamme : Polène, Céline, Jacquemus
      \item Mon terrain : \alert{Polène} --- déployé dans $>$20 pays, livré en continu
    \end{itemize}
    \column{.45\textwidth}
    \placeholder{3.2cm}{logo Polène / capture polene-paris.com}
  \end{columns}
\end{frame}

\begin{frame}{Un rôle qui s'est élargi en 3 temps}
  % 🎤 ~1 min 30. C'est LA slide qui annonce la largeur. Montrer la timeline, dire que je vais ZOOMER sur le 1er bloc.
  % 👉 A remplacer par une frise visuelle (TikZ ou image). Stub ci-dessous.
  \placeholder{3.4cm}{frise : Delivery QA  $\to$  Product / PO  $\to$  Discovery GTT (LNG)}
  \vskip 4pt
  \footnotesize
  \begin{columns}[T]
    \column{.33\textwidth}\textbf{Jan–Mar --- QA / Delivery}\\ Homologation, TNR, MEP multi-pays
    \column{.33\textwidth}\textbf{Avr–Mai --- Product}\\ Rédaction de tickets, audits UX/UI, appels d'offres
    \column{.33\textwidth}\textbf{Mai–Juin --- GTT / LNG}\\ Benchmark, discovery, ateliers
  \end{columns}
  \vskip 4pt
  % 🎤 « Je vais zoomer sur le premier bloc --- le projet le plus structurant --- puis montrer que la démarche se rejoue ailleurs. »
\end{frame}

\begin{frame}{Le fil rouge de cette soutenance}
  % 🎤 ~1 min 30. Poser le fil rouge transverse PUIS la problématique précise du projet phare.
  \begin{block}{Un même geste d'ingénieur, répété sur 3 missions}
    Transformer une réalité \alert{non structurée, non documentée, mouvante} en un livrable
    \alert{reproductible et transmissible} --- l'IA comme levier, mon jugement comme garde-fou.
  \end{block}
  \vskip 6pt
  \begin{alertblock}{Problématique du projet phare (QA)}
    Comment garantir la qualité du Delivery d'un site e-commerce multi-pays en exploitation,
    avec un dispositif opérationnel restreint et des cadences soutenues ?
  \end{alertblock}
\end{frame}

%========================================================================
\section{Le problème : une QA non reproductible}
% ⏱ ~4-5 min  |  Cœur du CADRAGE ingénieur. Le jury académique note ça.
%========================================================================

\begin{frame}{La QA dans le Delivery repose sur une seule condition}
  % 🎤 ~1 min. Les 3 missions QA, et la condition commune : la reproductibilité.
  \begin{columns}[T]
    \column{.6\textwidth}
    \begin{itemize}
      \item \textbf{Homologation} --- valider une nouvelle feature avant MEP
      \item \textbf{TNR} --- vérifier que l'existant n'est pas dégradé
      \item \textbf{QA au sens large} --- un niveau constant, des process reproductibles
    \end{itemize}
    \column{.4\textwidth}
    \begin{block}{Condition unique}
      Rejouer \alert{les mêmes vérifications}, sur \alert{les mêmes données}, dans \alert{les mêmes parcours}.
    \end{block}
  \end{columns}
\end{frame}

\begin{frame}{Trois défaillances empêchaient cette reproductibilité}
  % 🎤 ~2 min. Le cœur du diagnostic. Une défaillance = une cause racine, pas un symptôme.
  \begin{enumerate}
    \item \textbf{Tickets sous-spécifiés} --- critères d'acceptance flous, maquettes absentes
    \item \textbf{Documentation absente ou obsolète} --- seule la Home documentée, et dégradée
    \item \textbf{Pas de Golden Data} --- pré-prod partagée et modifiée quotidiennement $\Rightarrow$
          le référentiel testé change entre deux passes
  \end{enumerate}
  \vskip 6pt
  % 🎤 « Résultat : tests manuels exhaustifs avant chaque MEP, couverture inégale, régressions imprévisibles. »
\end{frame}

\begin{frame}{Première réponse, et pourquoi elle ne suffit pas}
  % 🎤 ~1 min 30. Montrer une réponse PRAGMATIQUE (QA Checks) + sa limite + la REFORMULATION (= maturité ingénieur).
  \begin{columns}[T]
    \column{.5\textwidth}
    \textbf{Les « QA Checks »}
    \begin{itemize}\footnotesize
      \item Section ajoutée au ticket avant dev
      \item Palliatif aux critères d'acceptance
      \item $\Rightarrow$ moins d'allers-retours dev/QA/PO
    \end{itemize}
    \column{.5\textwidth}
    \textbf{Mais limité par construction}
    \begin{itemize}\footnotesize
      \item Vit dans le périmètre d'un ticket
      \item Aucune vision globale du site
      \item Non automatisable
    \end{itemize}
  \end{columns}
  \vskip 8pt
  \begin{alertblock}{Reformulation (avec mon tuteur) $\to$ 2 décisions de conception indépendantes}
    Quelle \alert{stratégie de documentation} des TNR ? \quad Quel \alert{paradigme d'exécution} ?
  \end{alertblock}
\end{frame}

%========================================================================
\section{Décision 1 --- Documenter au bon niveau}
% ⏱ ~4 min  |  LE verrou méthodo. C'est ta meilleure réponse au « qu'avez-vous apporté ? »
%========================================================================

\begin{frame}{Quel niveau d'abstraction documenter ?}
  % 🎤 ~2 min. Le dilemme : trop fin = incassable à chaque MEP ; trop vague = inutile comme référence de test.
  \begin{columns}[T]
    \column{.5\textwidth}
    \textbf{L'expérience de la Home}
    \begin{itemize}\footnotesize
      \item Seule page pré-documentée : \alert{28 éléments} unitaires
      \item Chaque changement éditorial cassait des éléments
      \item Tout vérifier au moment d'une MEP : impossible dans le temps imparti
    \end{itemize}
    \column{.5\textwidth}
    \begin{block}{Deux extrêmes à éviter}
      \footnotesize
      Trop \alert{détaillé} (Cahier de recette) $\to$ non maintenable.\\[4pt]
      Trop \alert{vague} (scénarios courts) $\to$ pas une référence de test.
    \end{block}
  \end{columns}
\end{frame}

\begin{frame}{Le choix : documenter des blocs cohérents}
  % 🎤 ~2 min. Exemple Product Card. Et le couple « avant/après » de la Home comme ETALON de calibrage.
  \begin{itemize}
    \item Documenter \alert{composants \& fonctionnalités} comme des blocs (apparence, comportements, interactions)
    \item Ex. la \textbf{Product Card} = un bloc unique, et non 3 items séparés (add-to-cart, hover, variantes)
    \item Lisible, \alert{résistant aux changements éditoriaux}, suffisant comme référence TNR
  \end{itemize}
  \vskip 6pt
  \placeholder{2.6cm}{figure 4.2/4.4 du rapport : couple « avant (exhaustif) / après (bloc) »}
  % 🎤 « Ce couple avant/après est devenu mon étalon : il matérialise le niveau visé, et sert de référence à l'automatisation. »
\end{frame}

%========================================================================
\section{Construction --- Un pipeline de Skills}
% ⏱ ~4 min  |  Le « ce que j'ai construit ». Montrer l'ARCHITECTURE, pas le code.
%========================================================================

\begin{frame}{Automatiser la production de la documentation}
  % 🎤 ~2 min. Documenter 6 pages à la main = plusieurs semaines + incohérences.
  % Solution : une chaîne de 4 Skills Claude, calibrée sur l'étalon avant/après.
  \placeholder{3.0cm}{figure 2.1 du rapport : pipeline site-mapper $\to$ to-tnr $\to$ to-confluence / to-prompt}
  \vskip 4pt
  \footnotesize
  \begin{itemize}
    \item \texttt{site-mapper} --- cartographie exhaustive d'une page (via Playwright)
    \item \texttt{site-mapper-to-tnr} --- rationalise en fonctionnalités TNR au bon niveau
    \item \texttt{tnr-to-confluence} --- export Confluence (DOCX + CSV d'index)
    \item \texttt{tnr-to-prompt} --- transforme la doc en prompt exécutable par un agent
  \end{itemize}
\end{frame}

\begin{frame}{Deux partis pris d'ingénieur}
  % 🎤 ~2 min. Ce sont MES décisions d'architecture (pas l'IA). Bien les assumer.
  \begin{columns}[T]
    \column{.5\textwidth}
    \begin{block}{Séparation en 4 Skills}
      \footnotesize
      Responsabilité unique par étape $\Rightarrow$ on \alert{relance l'une sans les autres}.
      Générique : réapplicable à d'autres sites Shopify/SFCC en n'adaptant que l'entrée.
    \end{block}
    \column{.5\textwidth}
    \begin{alertblock}{Rétro-spec assumée}
      \footnotesize
      Faute de spec, on reconstruit depuis la prod.
      Limite connue : un bug en prod peut devenir une norme.
      Conséquence : tout écart pré-prod / prod sur un élément non modifié = bug.
    \end{alertblock}
  \end{columns}
  \vskip 6pt
  % 🎤 Maintenance : matrice par page dans Confluence, statut « Impacté MEP » avant chaque livraison.
\end{frame}

%========================================================================
\section{Décision 2 --- Le paradigme d'exécution}
% ⏱ ~5 min  |  Ton meilleur ARBITRAGE + ta contribution d'ingé la + nette (l'optim perf). Donne-lui de l'air.
%========================================================================

\begin{frame}{Une fois la doc produite : comment l'exécuter ?}
  % 🎤 ~1 min 30. Poser les deux paradigmes proprement.
  \begin{columns}[T]
    \column{.5\textwidth}
    \textbf{A --- Scénarios codés (outil dédié)}
    \begin{itemize}\footnotesize
      \item BrowserStack, Mr Suricate, Mabl, RainforestQA
      \item Déterministes une fois écrits
      \item Self-healing \alert{partiel} : vise les balises HTML, pas le contenu
    \end{itemize}
    \column{.5\textwidth}
    \textbf{B --- Exploration autonome par IA}
    \begin{itemize}\footnotesize
      \item L'agent déduit les scénarios de la doc $\to$ prompt
      \item Coût humain = \alert{lecture critique du rapport}
      \item Rapport vérifiable (captures + vidéos)
    \end{itemize}
  \end{columns}
  \vskip 4pt
  % 🎤 Réf clé pour le jury : Chevrot et al., « Are Autonomous Web Agents Good Testers? » (arXiv) --- sais le résumer !
\end{frame}

\begin{frame}{Comparaison sur 6 critères, et choix retenu}
  % 🎤 ~1 min 30. Le tableau est VALIDÉ dans le rapport --- ne pas le retoucher sur le fond.
  \scriptsize
  \begin{table}
    \begin{tabular}{lll}
      \toprule
      \textbf{Critère} & \textbf{A --- outil dédié} & \textbf{B --- exploration IA} \\
      \midrule
      Mise en place      & Élevé (1 j/scénario)        & Modéré (1 prompt) \\
      Maintenance        & Élevée (casse souvent)      & \alert{Faible (doc régénérée)} \\
      Réutilisabilité    & Limitée                     & \alert{Élevée (même pipeline)} \\
      Compétences        & Outil dédié / scripts       & Prompt en langage naturel \\
      Déterminisme       & \alert{Élevé}               & Modéré (variabilité LLM) \\
      Vérifiabilité      & Directe (assertions)        & Indirecte (rapport) \\
      \bottomrule
    \end{tabular}
  \end{table}
  \vskip 4pt
  \begin{block}{Choix : paradigme B + Antigravity (Google, vision Gemini)}
    \footnotesize Critère décisif = \alert{coût de maintenance} : aucune auto-réparation ne compense un
    référentiel éditorial mouvant ; le pipeline régénère la doc à la source.
  \end{block}
\end{frame}

\begin{frame}{Ma contribution la plus nette : l'optimisation perf}
  % 🎤 ~2 min. ★ SLIDE CLÉ. C'est CE que J'AI fait, pas l'IA. Quand on me demande « qu'avez-vous apporté ? » $\to$ ici.
  \begin{columns}[T]
    \column{.52\textwidth}
    \textbf{Constat}
    \begin{itemize}\footnotesize
      \item 1\textsuperscript{re} mesure : $\sim$\alert{1 h} par passe (Home, 9 fonctionnalités)
      \item Non tenable à l'échelle du site
      \item Cause : l'agent prenait lui-même captures \& vidéos en continu
    \end{itemize}
    \column{.48\textwidth}
    \textbf{Re-architecture}
    \begin{itemize}\footnotesize
      \item Sous-agent : exploration \alert{visuelle une fois}
      \item Puis \alert{Playwright} pilote captures
      \item $\to$ \alert{$\sim$20 min} par page
    \end{itemize}
  \end{columns}
  \vskip 6pt
  \begin{alertblock}{Le message à marteler}
    Le difficile n'était pas d'« utiliser des outils IA » mais de les \alert{orchestrer} :
    niveau d'abstraction, architecture du pipeline, et coût d'exécution maîtrisé.
  \end{alertblock}
\end{frame}

%========================================================================
\section{Résultats \& prise de recul}
% ⏱ ~4 min  |  Le jury académique CREUSE l'évaluation et les limites. Assume franchement.
%========================================================================

\begin{frame}{Résultats : du temps gagné au savoir transmis}
  % 🎤 ~1 min 30. Le gain de temps brut est dur à quantifier --- ASSUME-le et bascule sur le prouvable.
  \begin{itemize}
    \item \alert{6 pages prioritaires} documentées : Home, PLP, PDP, panier, checkout, compte
    \item Dispositif terrain : doc TNR (Confluence) + matrice de suivi + checklist avant MEP
  \end{itemize}
  \vskip 6pt
  \begin{block}{L'apport principal : le transfert de maîtrise}
    Avant : la connaissance des éléments à tester vivait dans \alert{quelques têtes}.
    Après : n'importe quel QA arrivant dispose d'un \alert{référentiel structuré}.
  \end{block}
\end{frame}

\begin{frame}{Limites --- assumées}
  % 🎤 ~1 min 30. Les dire AVANT que le jury ne les trouve = crédibilité.
  \begin{itemize}
    \item \textbf{Environnement non maîtrisé} --- pré-prod partagée ; pas de Golden Data ; duplication de store Shopify non native
    \item \textbf{Déterminisme} --- variabilité inhérente aux LLM (compensée par un rapport humainement vérifiable)
    \item \textbf{Couverture} --- Antigravity ne pilote que Chrome ; or les clients Polène sont surtout sur Safari mobile
    \item \textbf{Temps} --- une passe complète = plusieurs heures (OK avant MEP, contraignant à plus haute fréquence)
  \end{itemize}
\end{frame}

\begin{frame}{Impact : sobriété numérique vs. coût de l'IA}
  % 🎤 ~1 min. Volet RSE attendu par l'INSA. Honnête : un bénéfice indirect ET un coût direct.
  \begin{columns}[T]
    \column{.5\textwidth}
    \begin{examples}
      \footnotesize Mieux détecter les régressions de perf (poids, requêtes) $\Rightarrow$ énergie/bande passante
      économisées $\times$ le trafic.
    \end{examples}
    \column{.5\textwidth}
    \begin{alertblock}{}
      \footnotesize Mais chaque passe mobilise un LLM plusieurs minutes/page. À l'échelle,
      l'\alert{optimisation des tokens} deviendra prioritaire.
    \end{alertblock}
  \end{columns}
\end{frame}

%========================================================================
\section{La démarche se rejoue : GTT / LNG}
% ⏱ ~4 min  |  DOSE A REGLER avec Nathan : 1 slide (compact) ou 2 (développé).
% But : prouver que la démarche n'est pas liée aux outils QA --- même geste, autre domaine.
%========================================================================

\begin{frame}{Un tout autre contexte, la même démarche}
  % 🎤 ~2 min. GTT = brevet de cuves méthaniers, veut pivoter vers une solution digitale LNG.
  \begin{columns}[T]
    \column{.5\textwidth}
    \textbf{Le problème}
    \begin{itemize}\footnotesize
      \item Masse de docs, interviews, guides utilisateurs \alert{non structurés}
      \item Objectif : converger vers une solution digitale pertinente
    \end{itemize}
    \column{.5\textwidth}
    \textbf{Le même geste}
    \begin{itemize}\footnotesize
      \item Chaos $\to$ \alert{benchmark structuré} $\to$ synthèse $\to$ décision
      \item IA = levier (NotebookLM, Deep Research) ; jugement = garde-fou
    \end{itemize}
  \end{columns}
  \vskip 6pt
  \placeholder{2.2cm}{capture FigJam : tableau de synthèse benchmark / 7 catégories (optionnel)}
\end{frame}

\begin{frame}{Méthode benchmark $\to$ ateliers de convergence}
  % 🎤 ~2 min. OPTIONNELLE --- à couper si le timing serre. Montre la rigueur de discovery produit.
  \footnotesize
  \begin{enumerate}
    \item Liste de $\sim$20 solutions concurrentes
    \item \textbf{Gemini Deep Research} par outil $\to$ rapports, itérés pour la précision
    \item Consolidation dans \textbf{NotebookLM} (base interrogeable) + synthèse $\to$ \textbf{FigJam} (livrable client)
    \item Regroupement en \alert{7 catégories} ; sourcing systématique des citations
    \item Interviews $\to$ pain points $\to$ \alert{3 problématiques} crois\'ees avec le benchmark
    \item \textbf{Ateliers} GTT : tri / priorisation $\to$ cible (affréteurs / armateurs)
  \end{enumerate}
  \vskip 4pt
  % 🎤 Présenté en comité de pilotage GTT. Force de proposition + autonomie sur la méthode.
\end{frame}

%========================================================================
\section{Conclusion \& projet professionnel}
% ⏱ ~3 min  |  Refermer le fil rouge, dire la trajectoire, finir net.
%========================================================================

\begin{frame}{Un fil rouge tenu sur trois missions}
  % 🎤 ~1 min 30. Boucler : le même geste d'ingénieur, décliné 3 fois.
  \begin{columns}[T]
    \column{.33\textwidth}\textbf{QA}\\ \footnotesize doc + pipeline $\to$ TNR reproductibles, savoir transmis
    \column{.33\textwidth}\textbf{Product}\\ \footnotesize critères implicites $\to$ tickets explicites, structure
    \column{.33\textwidth}\textbf{GTT}\\ \footnotesize docs épars $\to$ benchmark structuré $\to$ décision
  \end{columns}
  \vskip 10pt
  \begin{block}{La constante}
    \alert{Chaos $\to$ artefact reproductible et transmissible}, avec l'IA comme levier
    et le jugement d'ingénieur comme garde-fou.
  \end{block}
\end{frame}

\begin{frame}{Ce que ce stage a clarifié}
  % 🎤 ~1 min 30. Bilan technique / humain / pro. Trajectoire QA -> Product. CDI.
  \begin{itemize}
    \item \textbf{Technique} --- concevoir et mettre en prod une méthodo de test IA sur un site réel
    \item \textbf{Humain} --- défendre des choix devant des décideurs ; doser confiance et humilité
    \item \textbf{Professionnel} --- trajectoire vers le \alert{Produit} / la gestion de projet (CDI signé)
  \end{itemize}
  \vskip 8pt
  % 🎤 Mention bonne pratique INSA : un LLM a appuyé reformulation/mise en forme/relecture du rapport ;
  %    le fond --- analyses, arbitrages, conclusions --- reste personnel.
  \centerline{\footnotesize\textcolor{gray}{Continuer à intégrer l'IA comme un levier au service du métier.}}
\end{frame}

\begin{frame}[allowframebreaks]{Références}
  \footnotesize
  \begin{thebibliography}{99}
    \bibitem[Antigravity]{antigravity}
      Google. \href{https://antigravity.google/docs/browser-subagent}{\emph{Antigravity, Browser Subagent}}, documentation officielle, 2026.
    \bibitem[Playwright]{playwright}
      Microsoft. \href{https://playwright.dev}{\emph{Playwright, End-to-end Testing for Modern Web Apps}}, 2026.
    \bibitem[BrowserStack]{browserstack}
      BrowserStack. \href{https://www.browserstack.com/events}{\emph{QA Leadership Summit 2026, QA Meets AI}}, Mars 2026.
    \bibitem[Chevrot, 2025]{chevrot}
      A. Chevrot \emph{et al.} (Avril 2025).
      \newblock \href{https://arxiv.org/abs/2504.01495}{\emph{Are Autonomous Web Agents Good Testers?}}
      \newblock arXiv:2504.01495 [cs.SE].
    \bibitem[LMArena]{lmarena}
      LMArena. \href{https://lmarena.ai/leaderboard}{\emph{Vision Arena Leaderboard}}, Mars 2026.
    \bibitem[GTT Digital]{gtt_digital}
      GTT. \href{https://www.gtt.fr/digital-solutions}{\emph{Digital Solutions, Smart Shipping}}, 2026.
    \bibitem[Deep Research]{gemini_dr}
      Google. \href{https://blog.google/products/gemini/google-gemini-deep-research/}{\emph{Try Deep Research in Gemini}}, 2024.
    \bibitem[NotebookLM]{notebooklm}
      Google Labs. \href{https://blog.google/technology/ai/notebooklm-google-ai/}{\emph{NotebookLM, an AI-first Notebook}}, 2023.
    \bibitem[Confluence]{confluence}
      Atlassian. \href{https://support.atlassian.com/confluence-cloud/docs/import-content-into-confluence-cloud/}{\emph{Import Content Into Confluence Cloud}}, 2026.
  \end{thebibliography}
\end{frame}


\begin{frame}
  \Huge{\centerline{Merci de votre attention}}
  \vskip 12pt
  \normalsize{\centerline{\textcolor{gray}{Des questions ?}}}
\end{frame}

\end{document}
